\documentclass{article}
% generated by Madoko, version 1.0.0-rc6
%mdk-data-line={1}


\usepackage[heading-base={2},section-num={False},bib-label={True}]{madoko2}


\begin{document}



{\mdfontfamily{SimSun}
%mdk-data-line={8}
\mdxtitleblockstart{}
%mdk-data-line={8}
\mdxtitle{\mdline{8}A Simple Mail Sever and Client}%mdk
\mdxauthorstart{}
%mdk-data-line={13}
\mdxauthorname{\mdline{13}翁健}%mdk
\mdxauthorend\mdtitleauthorrunning{}{}\mdxtitleblockend%mdk

%mdk-data-line={10}
\section{\mdline{10}1.\hspace*{0.5em}\mdline{10}前言}\label{section}%mdk%mdk

%mdk-data-line={11}
\noindent\mdline{11}SMTP（Simple Mail Transmission Protocol）和POP（Post Office Protocol）
虽然它们在安全性、命令集设定上存在着一些缺陷，但仍然是今天广泛使用的邮件收发协议。
按照Assignment 1的要求，我利用python的\mdline{13}\mdcode{socket}\mdline{13}接口实现了一个类\mdline{13}\mdcode{telnet}\mdline{13}的命令行工具，
作为SMTP和POP的客户端；同样的，也是利用\mdline{14}\mdcode{socket}\mdline{14}接口实现了与它交互的POP与SMTP服务器端。%mdk

%mdk-data-line={16}
%mdk-data-line={17}
\subsection{\mdline{17}1.1.\hspace*{0.5em}\mdline{17}Socket的选取}\label{sec-socket}%mdk%mdk%mdk

%mdk-data-line={19}
\noindent\mdline{19}因为邮件的收发是借由互联网完成的，所以需要选用\mdline{19}\mdcode{AF\_INET}\mdline{19}地址家族，
而同时SMTP和POP是两个基于连接的协议，所以应当选取\mdline{20}\mdcode{SOCK\_STREAM}\mdline{20}套接字。%mdk

%mdk-data-line={22}
\subsection{\mdline{22}1.2.\hspace*{0.5em}\mdline{22}消息的结尾}\label{section}%mdk%mdk

%mdk-data-line={23}
\noindent\mdline{23}因为邮箱协议要发送邮件，邮件有长有短，所以消息的结束以定长来结束是不可能的，为了解决这个问题，每条消息都以
\mdline{24}\textbackslash{}\mdline{24}r\mdline{24}\textbackslash{}\mdline{24}n（即回车和换行）为结尾。%mdk

%mdk-data-line={26}
\section{\mdline{26}2.\hspace*{0.5em}\mdline{26}客户端}\label{section}%mdk%mdk

%mdk-data-line={27}
\noindent\mdline{27}客户端中，我先利用\mdline{27}\emph{WireShark}\mdline{27}抓包观察\mdline{27}\mdcode{telnet}\mdline{27}产生的数据包，每次交互的命令都由命令加参数的形式组成。
然后我利用\mdline{28}\mdcode{socket}\mdline{28}连接发送消息的形式，也根据相同的命令制作相同的数据包，从而进行进程间通信。%mdk

%mdk-data-line={30}
\section{\mdline{30}3.\hspace*{0.5em}\mdline{30}SMTP服务器端}\label{sec-smtp}%mdk%mdk

%mdk-data-line={31}
\noindent\mdline{31}根据我对邮件收发整个流程的理解，其实每一层邮件代理服务器，相对于它的下一层，
它自己本身就是一个“客户端”，利用SMTP协议中的\mdline{32}\mdcode{telnet}\mdline{32}命令将这层邮件发给下一层。
所以我们并不需要为代理服务器，和最后的接受服务器区分开写不同的服务器端代码。%mdk

%mdk-data-line={35}
\mdline{35}为了实现SMTP最核心的邮件发送功能，我支持了SMTP的\mdline{35}\mdcode{HELO,MAIL,RCPT,DATA,QUIT,NOOP}\mdline{35}这几个核心命令。
和\mdline{36}\mdcode{telnet}\mdline{36}中类似，这几个命令都会在服务器完成响应之后，返一个回执给客户端。%mdk

%mdk-data-line={38}
\mdline{38}具体实现的详情，见下表：%mdk
\begin{mdtabular}{3}{\dimeval{(\linewidth)/3}}{1ex}%mdk
\begin{tabular}{lll}\midrule
\multicolumn{1}{c}{{\bfseries\mdline{41} SMTP命令}}&\multicolumn{1}{c}{{\bfseries\mdline{41}作用}}&\multicolumn{1}{c}{{\bfseries\mdline{41}备注}}\\

\midrule
\mdline{43}&\mdline{43} 进行握手&\mdline{43} 按照标准的SMTP都要传入一个hostname\\
\mdline{44}&\mdline{44} 确认连接&\mdline{44} 但是我发现不管传什么都是和\mdline{44}\mdcode{telnet}\mdline{44}开始\\
\mdline{45} HELO&\mdline{45}&\mdline{45} open的那个host进行握手，所以在我的实\\
\mdline{46}&\mdline{46}&\mdline{46} 现里面不传hostname我也进行握手确认连\\
\mdline{47}&\mdline{47}&\mdline{47} 接\\
\midrule
\mdline{49}&\mdline{49} 利用from:\mdline{49}\textless{}\mdline{49}xxx\mdline{49}\textgreater{}\mdline{49}&\mdline{49} 如果在完成了发件人的协商之后再次企图\\
\mdline{50} MAIL&\mdline{50} 参数进行发件人&\mdline{50} 调用这个参数则会引发错误\\
\mdline{51}&\mdline{51} 的协商&\mdline{51}\\
\midrule
\mdline{53}&\mdline{53} 利用to:\mdline{53}\textless{}\mdline{53}xxx\mdline{53}\textgreater{}\mdline{53}&\mdline{53} 如果没有完成发件人的协商就企图进行收\\
\mdline{54} RCPT&\mdline{54} 参数进行收件人&\mdline{54} 件人的协商，则会引发错误；收件人可以\\
\mdline{55}&\mdline{55} 的协商&\mdline{55} 有多个\\
\midrule
\mdline{57}&\mdline{57} 进行邮件正文&\mdline{57} 发件人和收件人两者缺一不可，否则会引\\
\mdline{58} DATA&\mdline{58} 内容的协商&\mdline{58} 发错误，以一个单个一行的（\mdline{58}\textquotedblleft{}\textbackslash{}r\textbackslash{}n.\textbackslash{}r\textbackslash{}n\textquotedblright{}\mdline{58}）\\
\mdline{59}&\mdline{59}&\mdline{59}为结束\\
\midrule
\mdline{61} QUIT&\mdline{61} 结束会话&\mdline{61} 同时退出客户端\\
\midrule
\mdline{63} NOOP&\mdline{63} 啥也不做&\mdline{63} 通常用于确认连接的完整性\\
\midrule
\end{tabular}\end{mdtabular}

%mdk-data-line={66}
\noindent\mdline{66}附表，依照SMTP标准，返回code以表示通信结果：%mdk
\begin{mdtabular}{2}{\dimeval{(\linewidth)/2}}{1ex}%mdk
\begin{tabular}{ll}\midrule
\multicolumn{1}{c}{{\bfseries\mdline{69}编号}}&\multicolumn{1}{c}{{\bfseries\mdline{69}原因}}\\

\midrule
\mdline{71}220&\mdline{71} 连接成功\\
\mdline{72}250&\mdline{72} 操作成功\\
\mdline{73}354&\mdline{73}开始协商\\
\mdline{74}&\mdline{74} 邮件正文\\
\mdline{75}503&\mdline{75} 其他错误\\
\mdline{76}501&\mdline{76}命令语法错误\\
\midrule
\end{tabular}\end{mdtabular}

%mdk-data-line={79}
\noindent\mdline{79}为了证明我的邮箱有收件功能，在协商完邮件正文之后，我就会在服务器本地开辟一个文件夹，
为\mdline{80}\mdcode{REPT}\mdline{80}的目标存下邮件，以便用于后面测试POP的功能。%mdk

%mdk-data-line={82}
\section{\mdline{82}4.\hspace*{0.5em}\mdline{82}POP服务器端}\label{sec-pop}%mdk%mdk

%mdk-data-line={87}
\begin{mdbmargintb}{4em}{}%mdk
\begin{mdflushright}%mdk
{\tiny\mdline{88}Created with~\href{https://www.madoko.net}{Madoko.net}.}%mdk
\end{mdflushright}%mdk
\end{mdbmargintb}%mdk
}%mdk


\end{document}
